% !TEX root = ../bb_AiRa_Kappaleet.tex

% ö  \equal  \%c3\%b6
% ä  \equal  \%C3\%A4

\xdef\sectiontitle{\Isectiontitle} %aiheuttama
\TOCrefs

%%%%%%%%%%%%%%%
\begin{frame}[label=1_1_a]%[label \equal 1_0_TOC1]
\frametitle{\texorpdfstring{\culine[\sectiontitleulinecolor]{\sectiontitle}}{\sectiontitle} }
\label{\TOCname}

\vspace{-0.5cm}

\begin{itemize}[leftmargin = 0.5cm,itemsep = 2mm]
    \item[1.1]<1-> \hyperlink{1_1_a}{\CDAlert[blue]{\IaSubsectiontitle}}
    \item[1.2]<1-> \hyperlink{1_2_a}{\IbSubsectiontitle}	
    \item[1.3]<1-> \hyperlink{1_3_a}{\IcSubsectiontitle}	
    \item[1.4]<1-> \hyperlink{1_4_a}{\IdSubsectiontitle}
    \item[1.5]<1-> \hyperlink{1_5_a}{\IeSubsectiontitle}
    \item[1.6]<1-> \hyperlink{1_6_a}{\IfSubsectiontitle}
    \item[1.7]<1-> \hyperlink{1_7_a}{\IgSubsectiontitle}
\end{itemize}

\vspace{-1cm}
\begin{itemize}[leftmargin = 8.5cm,itemsep = 2mm]
    \item[1.8]<1-> \hyperlink{1_8_a}{\IhSubsectiontitle}
	\item[1.9]<1-> \hyperlink{1_9_a}{\IiSubsectiontitle}
	\item[1.10]<1-> \hyperlink{1_10_a}{\IjSubsectiontitle}
	\item[1.11]<1-> \hyperlink{1_11_a}{\IkSubsectiontitle}
\end{itemize}

%\endofslide 
\end{frame}
%%%%%%%%%%%%%%%

\xdef\sectiontitle{\IaSubsectiontitle} 

%%%%%%%%%%%%%%%
\begin{frame}%[label \equal 1_1_a]
\frametitle{\texorpdfstring{\culine[\sectiontitleulinecolor]{\sectiontitle}}{\sectiontitle} }

\begin{itemize}[itemsep=1.7mm]
	\item Statistisen mekaniikan työkalujen avulla voidaan ymmärtää ja ennustaa \appear{\Alert[red]{systeemeitä, joiden hiukkaslukumäärä on valtava}}
%	Statistical mechanics makes predictions about properties and behaviour of \appear{\Alert[red]{systems whose number of particles is enormous}}
	
\item Mitä enemmän on hiukkasia/kappaleita\appear{, sitä enemmän on tuntemattomia suureita}
\item Suuri hiukkaslukumäärä parantaa \appear{systeemin käyttäytymisen ennustettavuutta:}\appear{ keskiarvojen} \appear{ja niiden kehityksen} \appear{edustavuus paranee} 

\item \CDAlert[red]{Termodynaaminen systeemi} on systeemi\appear{, jonka hiukkaslukumäärä on tarpeeksi suuri}\appear{, jotta ennusteista alkaa tulemaan tarkkoja}

\item \CDAlert[blue]{Statistinen mekaniikka}:
	\item[] \begin{itemize}[itemsep=2.5mm]
		\item klassinen
		\item ei ole 'modernia' fysiikkaa \appear{(ts. ei pohjaudu kvanttimekaniikkaan)}
		\item MUTTA \CDAlert[tunired]{sovelletaan monilla fysiikan osa-alueilla!}
	\end{itemize}
\end{itemize}

\endofslide 
%\endofsection
\end{frame}
%%%%%%%%%%%%%%%

%%%%%%%%%%%%%%%%
%\begin{frame}
%\frametitle{\texorpdfstring{\culine[\sectiontitleulinecolor]{\sectiontitle}}{\sectiontitle} }
%
%%\begin{minipage}{10.4cm}
%\begin{itemize}
%	\item Examine a two-sided room \appear{with $N$ point-like\\
%		 particles placed in the room} \appear{(see Fig.~1)}
%
%	\item For each particle\appear{, there are two alternatives:}\\ 
%		\appear{\Alert[red]{left}} \appear{or \Alert[blue]{right} } 
%	\item This is a thermodynamic \CDAlert[fuchsia]{'two-state system'} \\
%		\appear{(Recall coins or paramagnets in Thermal physics)}
%	\item Estimate in \CDAlert[fuchsia]{how many ways} we can place $N$ particles into the room \appear{so that $N_{R}$ number of particles are on the right-hand side?}
%	\item The answer is \CDAlert[fuchsia]{binomial coefficients}:	
%	\[
%\Omega \textrm{((} N, N_{R} \textrm{))}  \equal \appear{ W_{N_{R}}^{N} } \equal  \appear{\LP} 
%\begin{array}{l}
%N \\
%N_{R}
%\end{array} 
%\appear{\RP}   \equal \frac{N !}{N_{R} ! \textrm{((} N-N_{R} \textrm{))} !}
%\]
%\appear{giving us the '\Alert[tunired]{multiplicity}' or the ''\CDAlert[tunired]{number of ways to arrange...}"}
%
%\end{itemize}
%%\end{minipage}
%
%\xdef\loc{north east}
%\begin{tikzpicture}[remember picture,overlay] % anchor \equal south
%    \node (A) [yshift = 1*\figYshift,xshift = \figXshift, anchor = \loc] at (current page.\loc) {\includegraphics[page = 1,width = 4cm]{Figures_booklet/SOM_9_Statistical_mechanics_figures/SOM_Figure_1.png}}; %scale \equal \figscale. % 8cm max hei
%\end{tikzpicture}
%
%
%\endofslide 
%%\endofsection
%\end{frame}
%%%%%%%%%%%%%%%%


%%%%%%%%%%%%%%%
\begin{frame}
\frametitle{\texorpdfstring{\culine[\sectiontitleulinecolor]{\sectiontitle}}{\sectiontitle} }
\framesubtitle{$N = 4$}

\appear{}
%\begin{minipage}{10.4cm}
\begin{itemize}
	\item Kuinka neljä \appear{(\CDAlert[fuchsia]{erotettavissa olevaa})} \appear{hiukkasta voidaan järjestää kaksiosaiseen huoneeseen:}
\appear{	\xdef\loc{center}
\begin{tikzpicture}[remember picture,overlay] % anchor \equal south
    \node (A) [yshift = 1*\figYshift,xshift = -3.5cm, anchor = \loc] at (current page.\loc) {\includegraphics[page = 1,width = 5cm]{Figures_booklet/SOM_9_Statistical_mechanics_figures/SOM_Figure_2_a.png}}; %scale \equal \figscale. % 8cm max hei
\end{tikzpicture}
}
\appear{%\xdef\loc{north east}
\begin{tikzpicture}[remember picture,overlay] % anchor \equal south
    \node (A) [yshift = 4*\figXshift,xshift = 3.5cm, anchor = \loc] at (current page.\loc) {\includegraphics[page = 1,width = 5cm]{Figures_booklet/SOM_9_Statistical_mechanics_figures/SOM_Figure_2_b.png}}; %scale \equal \figscale. % 8cm max hei
\end{tikzpicture}}

\end{itemize}
%\end{minipage}



\endofslide 
%\endofsection
\end{frame}
%%%%%%%%%%%%%%%


%%%%%%%%%%%%%%%
\begin{frame}
\frametitle{\texorpdfstring{\culine[\sectiontitleulinecolor]{\sectiontitle}}{\sectiontitle} }
\framesubtitle{$N$ kasvaa}

\begin{minipage}{6.5cm}
\begin{itemize}[itemsep=1.9mm]
	\item Kun $N = 4$\appear{, on yhteensä 5 vaihto\-eh\-toista arvoa $N_{R}\;\appear{(\,= 0,1,2,3,4)$}}
	\item \appear{Kuvan~3a} \appear{jakaumassa $\Omega = \Omega \textrm{((} N_{R} \textrm{))}$ näkyy heikohko maksimi} $\appear{N_{R}} \equal \appear{2} \equal \appear{1 / 2 N}$ 
	
	\item Kun $N = 4$, jakauma on leveä\appear{, mutta jos $N$:ää kasvatetaan}\appear{, \Alert[red]{jakauma kaventuu}}
	
	\item Kaikissa tapauksissa maksimi saavutetaan kun $\appear{N_{R}}  \equal  \frac{1}{2} \appear{N}$
	\end{itemize}
\end{minipage}

\vspace{1.9mm}
\begin{itemize}[itemsep=1.9mm]
	
	
	\item Huomaa että suurilla ja samankokoisilla huoneilla \appear{ja $N$:n arvoilla}\appear{, molemmissa huoneissa on todennäköisesti lähes sama määrä hiukkasia}\appear{, vaikka yksittäisen hiukkasen paikasta ei tiedetä mitään}
\end{itemize}


\xdef\loc{north east}
\begin{tikzpicture}[remember picture,overlay] % anchor \equal south
    \node (A) [yshift = 1*\figYshift,xshift = \figXshift, anchor = \loc] at (current page.\loc) {\includegraphics[page = 1,width = 7cm]{Figures_booklet/SOM_9_Statistical_mechanics_figures/SOM_Figure_3.png}}; %scale \equal \figscale. % 8cm max hei
\end{tikzpicture}

\endofslide 
%\endofsection
\end{frame}
%%%%%%%%%%%%%%%


%%%%%%%%%%%%%%%
\begin{frame}
\frametitle{\texorpdfstring{\culine[\sectiontitleulinecolor]{\sectiontitle}}{\sectiontitle} }
\appear{}
\begin{itemize}
	\item Oletetaan hiukkaset 'ilma-molekyyleiksi'\appear{ ja huoneet samankokoisiksi,} \appear{ todennäköisintä on \Alert[red]{molemmissa huoneissa on lähes sama määrä molekyylejä}} 
	\item Tämä voidaan sanoa suurella suhteellisella tarkkuudella/luotettavuudella...
	\vspace{4.5cm}
	\item[] ...vaikkei yksittäisten molekyylien sijainteja tiedetäkään
\end{itemize}

\xdef\loc{center}
\begin{tikzpicture}[remember picture,overlay] % anchor \equal south
    \node (A) [yshift = 1.45*\figYshift,xshift = 0cm, anchor = \loc] at (current page.\loc) {\includegraphics[page =1,width =12cm]{Figures_booklet/SOM_9_Statistical_mechanics_figures/FYS_1627_SSP-4.pdf}}; %scale \equal \figscale. % 8cm max hei
\end{tikzpicture}
\endofslide 
%\endofsection
\end{frame}
%%%%%%%%%%%%%%%

%%%%%%%%%%%%%%%
\begin{frame}
\frametitle{\texorpdfstring{\culine[\sectiontitleulinecolor]{\sectiontitle}}{\sectiontitle} }
\framesubtitle{Mikro- ja makrotilat}

\appear{}
\begin{minipage}{8.5cm}
	\begin{itemize}[itemsep=1.7mm]
	\item Esimerkki selventää \CDAlert[red]{mikrotilan} \appear{ja \CDAlert[blue]{makrotilan} käsitteitä} 

\item Kuvassa 3d %
\xdef\loc{north east}%
\begin{tikzpicture}[remember picture,overlay] % anchor \equal south
    \node (A) [yshift =1*\figYshift,xshift=\figXshift, anchor=\loc] at (current page.\loc) {\includegraphics[page=1,width=4.8cm]{Figures_booklet/SOM_9_Statistical_mechanics_figures/SOM_Figure_3_d.png}}; %scale \equal \figscale. % 8cm max hei
\end{tikzpicture}%
 \appear{jakauman piikki nousee arvoon $10^{3 \cdot 10^{22}}$} \appear{, eli on olemassa $10^{3 \cdot 10^{22}}$ erilaista tapaa järjestää molekyylit huoneeseen}\appear{, jos molekyylit on jaettu kahteen yhtäsuureen osaan}\\
	 \end{itemize}
\end{minipage}
\vspace{-4.3mm}
\begin{itemize}[itemsep=1.7mm]
	\item[]  (vasempaan ja oikeaan) \appear{\textrm{((}joiden hiukkaslukumääriä kuvaa $ N_{\mathrm{R}} \textrm{))}$} 

\item Jokainen 'tapa' vastaa yhtä \CDAlert[red]{mikrotilaa}.
\appear{Jos tietäisimme jokaisen molekyylin tilan (sijainnin)}\appear{, voisimme määritellä jokaisen (mikro)tilan}

\item Kaikki $10^{3 \times 10^{22}}$ mikrotilaa \appear{kuuluvat ryhmään jossa molekyylit ovat jakautuneet tasaisesti kahteen yhtäsuureen puoliskoon.} \appear{Kaikilla $10^{3 \times 10^{22}}$ mikrotilalla $N_{R} \equal \frac{1}{2} \appear{N}$}

\item Tämä on yksi \CDAlert[blue]{makrotila}. \appear{Muilla makrotiloilla $N_{R}$:n arvo poikkeaa arvosta $N_{R} \equal \frac{1}{2} \appear{N}$}
\end{itemize}

\endofslide 
%\endofsection
\end{frame}
%%%%%%%%%%%%%%%

%%%%%%%%%%%%%%%
\begin{frame}
\frametitle{\texorpdfstring{\culine[\sectiontitleulinecolor]{\sectiontitle}}{\sectiontitle} }
\framesubtitle{Tasapainotila}

\seminormal
\appear{}
\begin{itemize}%[itemsep=2mm]
	\item Lähtemällä liikkeelle tilanteesta, jossa kaksiosaisessa huoneessa on huomattavan erisuuret määrät hiukkasia eri osissa\appear{, systeemi \CDAlert[tunired]{kehittyy kohti tilaa} \CDAlert[tunired]{jonka todennäköisyys on suurin}}\appear{, ts. kohti makrotilaa jonka mikrotilojen lukumäärä on suurin}\appear{, eli kohti tilannetta jossa molekyylejä on sama määrä molemmissa osissa huonetta}

\item Todennäköisintä makrotilaa kutsutaan \CDAlert[red]{tasapainotilaksi}

\item Termodynaaminen systeemi kehittyy spontaanisti kohti tasapainotilaa \appear{ja saavuttaa lopulta tasapainotilan}
\end{itemize}

% 6.5 cm  \equal  midway, 10 cm  \equal  3/4 
\vspace{2.5mm}
\begin{minipage}{8.5cm}
\begin{itemize}%[itemsep=2mm]
	\item Lähtiessä liikkeellä epätasapainotilasta\appear{, systeemin makro- ja mikrotilat muuttuvat}

\item \CDAlert[blue]{Tasapainossa} systeemin \Alert[tuniblue]{mikrotilat muuttuvat}\appear{, mutta ne kaikki kuuluvat yhteen \Alert[blue]{ja samaan makrotilaan}} \appear{(tai ovat sitä lähellä...)}\appear{, makrotilaa jolla on tietty  $ N_{\mathrm{R}}$}
\end{itemize}
\end{minipage}
\vspace{3mm}

\xdef\loc{south east}
\begin{tikzpicture}[remember picture,overlay] % anchor \equal south
    \node (A) [yshift=-0.25*\figYshift,xshift =\figXshift, anchor=\loc] at (current page.\loc) {\includegraphics[page=1,width=4.8cm]{Figures_booklet/SOM_9_Statistical_mechanics_figures/SOM_Figure_3_d.png}}; %scale \equal \figscale. % 8cm max hei
\end{tikzpicture}

%\endofslide 
\endofsection
\end{frame}
%%%%%%%%%%%%%%%



